\begin{abstract}
Rakshak-Edge classifies SMS messages, shared text, links and, on the Raspberry Pi~5 service, network flows as SCAM, UNCERTAIN or SAFE, and explains the evidence in plain language. It runs on Arm hardware only: a Raspberry Pi~5 service and an arm64-v8a Android app. The Pi service listens on loopback and runs in a network namespace with no route out, and the Android app's merged manifest removes the internet permission. Three experts feed one decision: a character CNN over the raw and canonical bytes of a message, a character CNN over URLs with an optional offline domain-intelligence scorer, and deterministic rules that also extract the evidence. The decision layer abstains when the detector is not confident and offers three modes, with quiet as the default. A Llama~3.2 1B Instruct model, quantised to INT4 and run through ExecuTorch, may write sentences about the evidence. A deterministic verifier checks each sentence against the evidence package and drops any sentence with an unsupported URL, number or instruction; a fixed template takes its place. The explainer never changes the verdict. On the Pi, quiet-mode verdicts in verdict-first mode had a median time to verdict of 4.11~ms\prose{} on a 40-message English run (p95 10.04~ms). We report campaign-disjoint test sets, independent sets marked as such, in-sample sets labelled, and failed pre-registered gates. The explainer has not yet run on the Android device.
\end{abstract}
